% Figure from MATH 452 notes, section 2. Compile with the notes preamble.
\begin{tikzpicture}[scale=1.1]
% the set E: a blob
\draw[thick, figB, fill=figB!12] plot[smooth cycle, tension=0.8] coordinates {(0,0) (2.2,-0.3) (3.2,0.8) (2.6,2.0) (1.0,2.3) (-0.4,1.2)};
\node[figB, font=\small] at (0.3,1.5) {$E$};
\node[black!70, font=\small] at (4.3,1.9) {$E^c$};
% interior point: a ball inside E
\draw[figG, thick, dashed, fill=figG!10] (1.3,1.1) circle (0.45);
\filldraw (1.3,1.1) circle (1.2pt);
\node[font=\scriptsize, anchor=north] at (1.3,0.62) {interior};
% boundary point (exactly on the curve: a control point of the cycle)
\draw[figR, thick, dashed] (3.2,0.8) circle (0.45);
\filldraw (3.2,0.8) circle (1.2pt);
\node[font=\scriptsize, anchor=south] at (3.2,1.27) {boundary};
% exterior point: a ball inside E^c
\draw[figG, thick, dashed] (4.6,0.55) circle (0.45);
\filldraw (4.6,0.55) circle (1.2pt);
\node[font=\scriptsize, anchor=north] at (4.6,0.07) {exterior};
\end{tikzpicture}
